\subsection{Evidence access can be a harder ceiling than reasoning scaffolds}

A final function-first search surfaced a particularly direct empirical comparator. Wang holds a production scientific-valuation agent's backbone fixed while ablating reasoning scaffolds/public tools and access to a curated proprietary evidence substrate \cite{wang2026evidence}. In that small stratified study, the structured reasoning layer improves calibration and audit discipline, while the much larger gain in factual coverage arrives only when the otherwise inaccessible evidence substrate is supplied. The study is domain-specific and its curated gold set is not an independent ground truth, so it does not establish a universal law of AI science. Its methodological implication is nevertheless important for the present architecture: an agent cannot reason its way to evidence that its acquisition boundary never exposes.

This narrows the role of \RAKL{} further. Evidence governance cannot compensate for an impoverished evidence universe; it can only make the limitation explicit, preserve acquisition provenance, and prevent missing evidence from being converted into unsupported authority. The intended advantage of the framework is therefore not that state discipline removes factual ceilings. It is that the ceiling, the search routes that produced it, and the claim types that remain blocked become first-class research state. A real superiority evaluation must manipulate reasoning architecture and evidence access separately rather than crediting one for the other's contribution.
